% Checkpoint35 proposed source; checkpoint34 sources remain unchanged.
\statement{A.112}{Proved}{The actual first congruence-central quotient is a specified product, and its odd coherent obstruction coordinates are binary}
Retain the original signed congruence map \(q:F\to\mathcal C\), put
\[
 \begin{gathered}
 N=\ker q,\quad N_1=\overline{[N,F]},\quad P=F^{(2)},\\
 M_2=\ker(P\to\mathcal C_{(2)}),\quad
 M_{21}=\overline{[M_2,P]},\quad P_*=P/M_{21}.
 \end{gathered}
\]
and write \(G_p=\mathrm{SL}_2(\mathbb Z_p)\).
Let \(\chi:G_3\to C_3\) be its actual abelianization character,
normalized by \(\chi(q_3(x))=\chi(q_3(y))=1\).
There is an isomorphism of actual compact groups
\[
 E_1=F/N_1\ \cong\
 P_*\times\bigl(G_3\times_{C_3}\mathbb Z_3^2\bigr)
       \times\prod_{p\geq5}\bigl(G_p\times\mathbb Z_p^2\bigr).
 \tag{A.112.1}
\]
The fiber product uses \(\chi\) and \((a,b)\mapsto a+b\bmod3\).
All maps in this identification are the original coordinate maps, and
\[
 E_1'=P_*'\times G_3'\times\prod_{p\geq5}G_p.
 \tag{A.112.2}
\]
Put \(J=\left(\begin{smallmatrix}0&-1\\1&-1\end{smallmatrix}\right)\).
For \(e\in\{1,-1\}\), set \(\varepsilon=(1-e)/2\),
\[
 R_e=y^\varepsilon r,\qquad \rho_e=\operatorname{Ad}(R_e),\qquad
 J_e=\operatorname{diag}(1,e)J\operatorname{diag}(1,e)^{-1}.
\]
Then
\[
 \begin{split}
 T_e=C_{E_1}(R_e)
   &=(P_*')^{\rho_e}\times T_{3,e}\times\prod_{p\geq5}T_{p,e},\\
 T_{3,e}&=C_{G_3'}(J_e),\qquad T_{p,e}=C_{G_p}(J_e)\quad(p\geq5).
 \end{split}
 \tag{A.112.3}
\]
For the actual full-GT negative involution and the ORIGINAL negative
frame \(J_e^\alpha\) of A.111, let
\(\kappa_e(t)=(J_e^\alpha)^{-1}\alpha(t)J_e^\alpha\).
This action is inversion on every odd factor in (A.112.3), and
\[
 H^1(C_2,T_e;\kappa_e)\ \cong\
 H^1(C_2,(P_*')^{\rho_e};\kappa_e)\times
                  \prod_{p\ {\rm odd}}C_2
 \tag{A.112.4}
\]
as pointed compact spaces. Its three-adic coordinate is the actual
scalar reduction \(q_3\). For the specified coherent class of A.111,
that coordinate is
\(\epsilon_A(d\bmod3)I=(-1)^{\Delta(\alpha)}I\).
No value of this coordinate, of the other odd coordinates, or of the
dyadic class is asserted.

\begin{proof}
\emph{A.112.1. The exact image of the original coordinates.}
A.106 gives
\[
 N_1=N\cap F'\cap\pi_2^{-1}(M_{21}).
\]
Thus the original congruence, abelian, and \(P_*\) coordinates embed
\(E_1\) into \(\mathcal C\times F^{\rm ab}\times P_*\).
The dyadic congruence and abelian coordinates are determined by \(P_*\).
Writing \(A_o=\prod_{p\ {\rm odd}}\mathbb Z_p^2\) and
\(G_o=\prod_{p\ {\rm odd}}G_p\), we obtain an actual embedding
\[
                  E_1\longrightarrow (P_*\times A_o)\times G_o.
 \tag{A.112.5}
\]
The joint map to \(P_*\times A_o\) is onto: its closed image is
subdirect, and a common quotient of a pro-two group and a pro-odd
group is trivial. The map onto \(G_o\) is onto as well.

We spell out the compact Goursat argument. If \(I\leq U\times V\)
is closed and projects onto both factors, set
\(U_0=\{u:(u,1)\in I\}\) and \(V_0=\{v:(1,v)\in I\}\).
They are closed normal subgroups, and a pair in \(I\) identifies
\(U/U_0\) continuously with \(V/V_0\). Compactness supplies the quotient
topology. The subgroup \(I\) is exactly the inverse image of this
isomorphism's graph.

The primary abelianization calculation is
\[
              G_3^{\rm ab}=C_3,\qquad G_p^{\rm ab}=1\quad(p\geq5),
 \tag{A.112.6}
\]
in the precise scope of Beyl's Proposition~5.7 used in
A.106--A.108.\citep[Proposition~5.7]{A_Beyl1986}
It implies that the maximal pronilpotent quotient of \(G_o\)
is exactly the displayed \(C_3\). Indeed a finite \(p\)-group with
cyclic abelianization is cyclic: a lift of an abelianization
generator, if contained in a proper maximal subgroup, contradicts
the fact that this subgroup contains the commutator subgroup.
A \(p\)-group with zero abelianization is therefore trivial.
A continuous map from \(G_o\) to a finite \(p\)-group depends on
only finitely many prime coordinates. Every \(G_\ell\), \(\ell\geq5\),
has trivial image; the \(G_3\) image is trivial unless \(p=3\),
and then is a quotient of its actual \(C_3\).
Applying the Sylow projections proves the assertion for every finite
nilpotent quotient and hence every pronilpotent quotient.

The group \(P_*\times A_o\) is pronilpotent. Its common quotient with
\(G_o\) in (A.112.5) is consequently a quotient of \(C_3\), so the
actual image has index at most three in the product. On the other
hand its original generators satisfy
\[
             \chi(q_3(g))=\exp_x(g)+\exp_y(g)\pmod3.
 \tag{A.112.7}
\]
This confines the image to the indicated index-three fiber product.
The two index bounds force equality. Rearranging its coordinates
proves (A.112.1).

\emph{A.112.2. Derived groups and based cyclic actions.}
The derived group of \(G_3\times_{C_3}\mathbb Z_3^2\) is
\(G_3'\times\{0\}\): lift any pair from its surjection onto \(G_3\)
and take the commutator. Closure gives all of \(G_3'\).
For an infinite direct product, the closed derived subgroup contains
each supported derived factor and hence their dense finite products.
This proves (A.112.2).

On \(F'\) the extra signed character in the original cyclic
congruence formula vanishes. Thus \(\rho=\operatorname{Ad}(r)\)
acts on every odd factor as \(\operatorname{Ad}(J)\).
The dyadic action is also defined on \(P_*\): the modulo-four
calculation in A.107 shows that an element of \(M_2\) has both
exponent parities zero, so the signed cyclic formula preserves
\(M_2\); order three gives equality. It then preserves
\(M_{21}=\overline{[M_2,P]}\).
The actual abelian cyclic matrix has no fixed vector because
\(\det(\rho-I)=3\) and multiplication by three on
\(\widehat{\mathbb Z}^2\) is injective.
Hence every pure centralizer of \(R_e\) lies in \(E_1'\).
The exact integral identity
\[
       J_e=(-1)^\varepsilon Y^\varepsilon J,\qquad
       Y=\begin{pmatrix}1&0\\-2&1\end{pmatrix},
\]
now proves (A.112.3), including its actual dyadic factor.

\emph{A.112.3. The odd factors and their square quotients.}
The commutation equations with \(J_e\), which has a cyclic-vector
basis, give its full matrix centralizer algebra
\(\mathbb Z_p[J_e]\). Determinant is the quadratic norm.
Its norm-one unit group is therefore \(C_{G_p}(J_e)\), an abelian
compact group. For \(p\geq5\), its quadratic discriminant \(-3\)
is a unit. The residue norm-one group is cyclic of order \(p-1\)
in the split case and \(p+1\) in the nonsplit case.
Reduction is onto: lift a residue element and correct its
determinant in \(1+p\mathbb Z_p\) by a scalar square root.
The kernel is abelian pro-\(p\), where squaring is invertible.
Consequently
\[
                         T_{p,e}/T_{p,e}^2=C_2\quad(p\geq5).
 \tag{A.112.8}
\]
At three, the already bound finite centralizer is
\(\langle J_e,-I\rangle\) of order six. Its three-part has nonzero
actual \(\chi\)-value, whereas \(G_3'=\ker\chi\).
Thus \(T_{3,e}\) reduces exactly onto \(\{I,-I\}\), with abelian
pro-three kernel \(U_{3,e}\). The central \(-I\) splits it:
\[
 T_{3,e}=\{I,-I\}\times U_{3,e},\qquad
                         T_{3,e}/T_{3,e}^2=\{I,-I\}.
 \tag{A.112.9}
\]
All square images are closed, being images of compact groups.
The finite centralizer input is the same one certified in
A.105/A.111; no new finite search is invoked.

\emph{A.112.4. The coherent action and the exact product of spaces.}
The original relation
\(\alpha(R_e)=J_e^\alpha R_e^{-1}(J_e^\alpha)^{-1}\)
makes \(\kappa_e\) conjugate \(\rho_e\) to \(\rho_e^{-1}\).
The actual dyadic and odd coordinate maps show that it acts
componentwise in (A.112.3).
At an odd prime write its matrix action as conjugation by \(L_p\).
Then conjugation by \(L_pJ_eL_p^{-1}\) agrees on the derived factor
with conjugation by \(J_e^{-1}\).
The matrix centralizer of that derived factor is scalar.
For \(p\geq5\) use its two elementary unipotents; at three use
\(X^3=I+6E_{12}\) and \(Y^3=I-6E_{21}\) in \(G_3'\).
The integral domains \(\mathbb Z_p\) make these commutation
equations scalar even at three. Hence
\(L_pJ_eL_p^{-1}=\zeta_pJ_e^{-1}\).
Both cyclic matrices have trace \(-1\), forcing \(\zeta_p=1\).
It follows that \(\kappa_e\) is the quadratic conjugation on
\(\mathbb Z_p[J_e]\), and is inversion on norm-one units.

For an abelian compact group acted on by inversion, all elements
are cocycles and the coboundaries are the squares.
Cocycles and correcting conjugators for a product action are
componentwise. Compact orbit quotients are Hausdorff, so these
identifications are homeomorphisms. Equations (A.112.8)--(A.112.9)
give (A.112.4). A.111's same signed map identifies its three-adic
coordinate with the specified original CM sign. The decomposition
does not evaluate that sign.
\end{proof}

\begingroup\footnotesize\raggedright
Proof record: the complete checkpoint34 source
\path{work_DE/actual_central_quotient_product_and_odd_obstruction34.md}
and both entire-source independent reviews are retained.
This formal source uses the actual coordinate image throughout;
it does not replace the original negative frame by an arbitrary
quotient frame.\par\endgroup
